% SOURCE_AUTHORITY: sga5_fr_workpass.tex, lines 12117--12131 (LF-normalized unit).
% SOURCE_SHA256: 791F4EFFC5E02832D5D77ED03518C8156D6F07E4C8238B03545DB93D883FBB28
Se $A$ é um anel comutativo, e se $X^\bullet$, $Y^\bullet$ são dois complexos perfeitos, então seu produto tensorial
\[
X^\bullet\overset{\mathbf L}{\otimes}_A Y^\bullet
\]
na categoria derivada continua a ser um complexo perfeito. Obtém-se um bifuntor
\[
D(A)_{\parf}\times D(A)_{\parf}\longrightarrow D(A)_{\parf}
\]
exato em cada variável. Por conseguinte, define-se um produto em $K^{\bullet}(A)$, e é fácil ver que $K^{\bullet}(A)$ fica assim dotado de uma estrutura de anel comutativo com elemento unidade. De fato, $K^{\bullet}(A)$ pode ser dotado de uma estrutura de $\lambda$-anel (vejam-se os detalhes no Seminário sobre Riemann--Roch que se segue ao presente\footnote{SGA 6 V 2.}).

O produto tensorial
\[
X^\bullet\overset{\mathbf L}{\otimes}_A Y^\bullet
\]
de um complexo perfeito $X^\bullet$ e um complexo pseudocoerente limitado $Y^\bullet$ é um complexo pseudocoerente limitado (sendo $A$ sempre um anel comutativo). Por conseguinte, pode-se introduzir em $K_\bullet(A)$ uma estrutura de $K^{\bullet}(A)$-módulo.
